
\documentclass[11pt]{article}

\usepackage[T1]{fontenc}
\usepackage{graphicx,wrapfig,lipsum}
\usepackage{hyperref}
\usepackage{xcolor}
\hypersetup{colorlinks,
  linkcolor={red!50!black},
  citecolor={blue!50!black},
  urlcolor={blue!80!black}}
\usepackage{enumerate}
\usepackage{amsmath}
\usepackage{amssymb}
\usepackage{mathtools}
\usepackage{geometry}
 \geometry{
 a4paper,
 total={170mm,257mm},
 left=20mm,
 top=20mm,
}
\renewcommand{\thefootnote}{$\dagger$} 

\title{Scipy: Large-scale Constrained Optimization}
\author{Ant\^{o}nio Horta Ribeiro}
\date{}

\begin{document}

\maketitle

\section*{Sub-organization: Scipy}

\section*{Student Info}

\begin{itemize}
\item
  Name: Ant\^{o}nio Horta Ribeiro
\item
  GitHub: \includegraphics[height=\baselineskip]{GH.png} \href{https://github.com/antonior92}{antonior92}
\item
  Email: \href{mailto:antonior92@gmail.com}{antonior92@gmail.com}
\item
  Website: \href{https://antonior92.github.io}{antonior92.github.io}
\item
  CV: \href{https://www.dropbox.com/s/3qgyuasajp9pjpy/cv_8.pdf?dl=0}{https://www.dropbox.com/s/3qgyuasajp9pjpy/cv\_8.pdf?dl=0}
\item
  Tel.: +55 (31) 99741-9213
\item
  Timezone: GMT-3
\item
  GSoC Blog RSS Feed URL:  \href{http://hortaribeiro.wordpress.com/?feed=rss}{http://hortaribeiro.wordpress.com/?feed=rss}
\end{itemize}


\section*{University Info}

\begin{itemize}
\item
  University Name: Federal University of Minas Gerais (Brazil)
\item
  Major: Electrical Engineering
\item
  Current Year: 2nd year
\item
  Degree: Master
\end{itemize}


\section*{Code Sample}

I have contributed to Scipy optimization and signal processing packages in the following
pull requests. 

\begin{itemize}
  \item
    ENH: signal: added irrnotch and iirpeak functions. \#6404 --- \\
    \href{https://github.com/scipy/scipy/pull/6404}{https://github.com/scipy/scipy/pull/6404}
    (merged)
  \item
    ENH: optimize: added iterative trustregion. \#6919 --- \\
    \href{https://github.com/scipy/scipy/pull/6919}{https://github.com/scipy/scipy/pull/6919}
    (merged)
  \item
    ENH: optimize: changes to make BFGS implementation more efficient. \#7165 --- \\
    \href{https://github.com/scipy/scipy/pull/7165}{https://github.com/scipy/scipy/pull/7165}
    (merged)
\end{itemize}

\noindent
Both pull requests: \#6919 and  \#7165, showcase my optimization skills.

\newpage
\section*{Project Info}

\subsection*{Title: \\``\textit{Scipy: Large-scale Constrained Optimization}''}

\subsection*{Abstract:}


\textbf{\noindent
  The project consists of implementing in Scipy an interior-point
  method for nonlinear problems. The main goal is to implement
  a constrained optimization algorithm able to deal with a large
  (and possibly sparse) problems for which the constrained optimization methods currently
  implemented in Scipy (namely \texttt{SLSQP} and \texttt{COBYLA}) are
  largely inappropriate to deal with. Implement benchmark problems and
  integrate quasi-Newton (namely BFGS, SR1 and L-BFGS)
  and finite differences (using graph coloring schemes to deal with sparse structures)
  approximations to the method are also part of the project.}

\subsection*{Description:}

This project intends to deal with the problem of minimizing a function
subject to constraints on its variables, mathematically formulated as
the following nonlinear programming (NLP) problem:
%
\begin{eqnarray}
  \label{eq:NLP}
   \min_x && f(x)\\\nonumber
  \text{subject to } && c_E(x) = 0 \\\nonumber
  && c_I(x) \ge 0,
\end{eqnarray}
%
where $x\in \mathbb{R}^n$ is a vector of unknowns and the objective function $f$,
the equality constraints $c_E$ and the inequality constraints $c_I$ are twice continuous
differentiable functions. In science, engineering and finance
many applications can be formulated as the above optimization problem.

The solvers for problem (\ref{eq:NLP}) can be
classified into three groups, described next (together with a list
of the main optimization softwares that uses each approach): 
%
\begin{enumerate}[(i)]
\item Augmented Lagrangian methods (MINOS~\cite{murtagh1983minos}
  and LANCELOT~\cite{conn1992lancelot});
\item SQP (Sequential Quadratic Programming) methods (SNOPT~\cite{gill2005snopt}, 
FILTERSQP~\cite{fletcher1998user} and \\KNITRO/ACTIVE~\cite{byrd2006knitro});
\item Interior-point methods (KNITRO/DIRECT,
  KNITRO/CG~\cite{byrd2006knitro}, IPOPT~\cite{wachter2006implementation},
  LOQO~\cite{vanderbei1999loqo} and \\BARNLP~\cite{barnlp2000boing}).
\end{enumerate}
%
While the first widely available softwares for large-scale optimization
(MINOS and LANCELOT) adopted augmented Lagrangian methods, more
modern softwares usually adopt interior-point or SQP formulations. The existing
benchmarks~\cite{morales2003assessing, benson2003comparative, dolan2006optimality}
don't offer a single approach as the best for all cases and SQP and interior-point
formulations often appears as competing state-of-the-art approaches with
complementary strengths. According to Nocedal and Wrigth~\cite[p.560 and p.592]{nocedal2006numerical},
SQP methods are the most efficient if the number of constraints is nearly as large
as the number of variables and are more robust to badly scaled problems,
while interior point methods show their strength in large-scale applications,
where they often (but not always) outperform SQP methods. The rapid pace of software
development, however,  makes it difficult to arrive at concrete conclusions about
the classes of problems each approach is best suited for.

Since Scipy already have a SQP method (\texttt{SLSQP}) available, an interior
point method seems the most sensible choice for this project. Furthermore,
interior-points methods often have a good performance in large-scale
applications~\cite[p.592]{nocedal2006numerical}, while maintaining a
competitive performance for medium/small problems as well~\cite{morales2003assessing}.

The algorithm to be implemented in this project is the interior-point
trust-region method described in~\cite{byrd1999interior}.
This implementation has been chosen because of its ability to deal with large
problems with, possibly, sparse structure. The trust-region approach allows
great freedom in the choice of the Hessian and provides a mechanism for coping
with Jacobian and Hessian singularities. Furthermore, it is a widely tested method,
being the algorithm implemented in the commercial package KNITRO/CG~\cite{byrd2006knitro};
There are comprehensive references about the algorithm and its variations
(namely~\cite{byrd1999interior, byrd2000trust, nocedal2006numerical}); And,
while the algorithm is not very new (it was first described in 1999),
it is still object of active search (improvements in the infeasibility detection
of this algorithm have just been published in 2014~\cite{nocedal2014interior}).


\subsection*{Schedule and Deliverables:}

Project milestones are described next. 

\subsubsection*{Further Literature Review --- Before the official code begins}

\begin{itemize}
\item
  The success of this project is largely depending on a deep knowledge about the methods being
  implemented and about optimization theory in general. So I intend is to read as much as possible
  about the theme before the official code period begins so things go as smoothly
  as possible during the implementation.
\end{itemize}

\subsubsection*{Prepare for  the Project and Interact with Mentors --- Community bounding period}

\begin{itemize}
\item
  I already have some familiarity with \texttt{scipy.optimize} internal structure
  from previous contributions. That way, my main concern during the bounding period is
  to interact with my mentors and figure out small details about the algorithms
  I intend to implement.
\end{itemize}


\subsubsection*{Equality-constrained Quadratic Programming --- 1  week}
\begin{itemize}
\item
  The interior point method~\cite{byrd1999interior} need to solve, in one of its sub-steps,
  an equality-constrained quadratic programming (EQP) problem:
  %
  \begin{eqnarray}
    \label{eq:eqp}
    \min_x && \tfrac{1}{2}x^THx+g^Tx\\\nonumber
    \text{subject to } &&  A x = b.
  \end{eqnarray}
  %
  Nocedal and Wright~\cite[Sec. 16.1 to 16.3]{nocedal2006numerical} describe
  several  methods for solving this problem.
  Among them I intend to implement two: the so-called projected conjugate
  gradient (CG) method  and the null-space method. Both of them can deal with indefinite
  matrix $H$ and have complementary strengths: the projected CG method is an iterative
  method, it have low memory-storage requirements and may provide an approximate
  solutions to the problem within few iterations, being a good choice
  for large-scale problems; and, the null-space method is a direct factorization method,
  it provide accurate solutions and may be a better choice for small/medium-sized
  problems. Both methods can be adapted to deal with sparse structures.
\end{itemize}


\subsubsection*{Trust-region Interior Point Method --- 3 to 5 weeks}

\begin{itemize}
\item
  At this stage of the project, I intend to implement, document and test the trust-region
  interior point method described in~\cite{byrd1999interior}. The general idea of
  the method is explained next.

  Consider the barrier problem:
  \begin{eqnarray}
    \label{eq:barrier_problem}
       \min_x && f(x) - \mu \sum \log s_i\\\nonumber
       \text{subject to } && c_E(x) = 0 \\\nonumber
       && c_I(x)  = s.
  \end{eqnarray}
  One need not to include inequality constraints because the barrier term $ - \mu \sum \log s_i$
  prevents the components of $s$ from becoming close to zero (Because $(-\log s_i) \rightarrow \infty$
  when $s_i \rightarrow 0$) and guarantee the inequality constraints are being satisfied.
  Interior-point methods finds approximate solutions to the above barrier problem~(\ref{eq:barrier_problem})
  for a sequence of positive barrier parameters $\{\mu_k\}$ that converges to zero.
  And, as $\mu \rightarrow 0$ the solution of the barrier problem approach the
  solution of the NLP problem~(\ref{eq:NLP}).
  
  The interior point method to be implemented solve the barrier problem~(\ref{eq:barrier_problem})
  by sequentially solving a series of EQP problems subjected to  trust-region constraints.
  The EQP solvers implemented during the previous weeks are needed as sub-steps of the method.
\end{itemize}

\subsubsection*{Benchmarks --- 1 or 2 weeks}

\begin{itemize}
\item

  It is expected that small tests and problems will be tried along the development of the algorithm.
  At this stage, however, I intend to benchmark the algorithm in a more rigorous manner,
  implementing a subset of the problems from  CUTEst~\cite{gould2015cutest}
  in Python and testing the implemented method on it.  Furthermore, I intend to compare the method
  with other available optimization software. For instance, IPOPT~\cite{wachter2006implementation}
  package already have a wrapper for Python, so it should be really easy to compare with it.
\end{itemize}

\subsubsection*{Preconditioners --- 1 week}
\begin{itemize}
\item
  Preconditioning (described in~\cite[pp. 118-120]{nocedal2006numerical})
  is a method of improving the convergence of CG implementations.
  Currently there are not any preconditioner implemented in Scipy and
  I intend to implement and test preconditioners for the projected CG method.
  This will improve the performance of the interior-point methods,
  since the projected CG is one of its sub-steps.
\item
  Both \texttt{trust-ncg} and \texttt{Newton-CG} could benefit from the
  implemented preconditioner, therefore I intend to include preconditioners options on both.
\end{itemize}

\subsubsection*{Finite-difference Approximations --- 2 weeks}


\begin{itemize}
\item
  The implemented algorithm require the first and second derivatives of functions and constraints.
  What may not easy to compute in some cases, so it is very useful to have some automatic way to
  estimate these derivatives. At this stage of the project I intend to implement finite-difference methods to
  estimate derivatives.

  Consider the definition of partial derivative:
  \begin{equation*}
    \frac{\partial f}{\partial x_i} = \lim_{\epsilon \rightarrow 0} \frac{f(x+\epsilon e_i) - f(x)}{\epsilon}.
  \end{equation*}
  The basic motivation of finite-difference approximations is to consider that:
  \begin{equation*}
    \frac{\partial f}{\partial x_i} \approx  \frac{f(x+\epsilon e_i) - f(x)}{\epsilon},
  \end{equation*}
  for a sufficiently small $\epsilon$. The above is the so-called forward-difference approximation,
  an alternative formulation is, for instance, the central-difference approximation.

  Given a scalar function $f: \mathbb{R}^n \rightarrow \mathbb{R}$ one can estimate its
  gradient $\nabla f$ and its Hessian $\nabla^2 f$ using finite-differences.
  A vector function $c: \mathbb{R}^n \rightarrow \mathbb{R}^m$ may also have its Jacobian $J$
  matrix estimated using finite differences. Some of the details of the implementation
  (e.g. the optimal choice of $\epsilon$) are covered in~\cite[Sec. 8.1]{nocedal2006numerical}.
\item
  The Hessian $\nabla^2 f$ and the Jacobian $J$ may have a known sparse structure. For this situations,
  the estimation may be performed within fewer functions evaluations using graph coloring
  schemes~\cite{coleman1983estimation, coleman1984estimation}. I intend to implement and
  test these schemes.

  Some work to approximate sparse Jacobians have already been done
  (functions \texttt{approx\_derivative} and \texttt{check\_derivative}) and can be reused.
  Sparse Hessians approximation algorithms will have to be implemented from scratch, since they are different
  from sparse Jacobian algorithms because of the Hessian symmetry.
  
  Its worth notice that this graph coloring schemes works almost
  without modifications for other derivatives calculation approaches.
  That way, if more advanced methods for calculating the derivatives
  (e.g. automatic differentiation~\cite[Sec. 8.2]{nocedal2006numerical})
  are ever implemented in Scipy they could reuse this functions to deal
  with sparse structures.
\end{itemize}

\subsubsection*{Quasi-Newton Hessian Approximations\footnotemark --- 1 or 2 weeks}

\footnotetext{This is a low priority item.  It implementation is dependent
  on the pace of the project and may be simplified or suppressed if needed.}


\begin{itemize}
\item
  It is useful to be able to use quasi-Newton Hessian approximation
  methods together with the interior point method. I intend to implement,
  integrate and test three quasi-Newton approximation methods
  for the use with the implemented interior point method:
  BFGS, SR1 and L-BFGS~\cite[Cap. 6 \& 7]{nocedal2006numerical}

  Scipy already have a \texttt{BFGS} implementation that could be reused,
  but SR1 and L-BFGS would have to be implemented from scratch (Scipy have
  a \texttt{L-BFGS-B} method available but it would not be of much use
  because it is only a wrapper for a FORTRAN function)
\end{itemize}



\subsubsection*{Tie Up the Loose Ends --- Last week}

\begin{itemize}
\item
  Finish to document, debug and optimize the code.
\end{itemize}



\section*{Other Commitments}

I will attend to \href{https://www.ifac2017.org}{IFAC world congress} from 9 to 14 of July.

\bibliography{GSoC2017}
\bibliographystyle{ieeetr}

\end{document}
